
% ----------------------------------------------------------------------------
\subsection{Dataset and Preprocessing}
% ----------------------------------------------------------------------------

We use the Quora Question Pairs dataset, a binary duplicate-detection benchmark
of \textbf{404{,}290} question pairs. Each example contains two natural-language
questions and a label indicating whether they are semantically equivalent. The
task is a strong testbed for model comparison because it mixes shallow lexical
overlap with paraphrase matching, so a model must reason about both surface form
and meaning.

Two properties of the data drive every later design choice. First, the classes
are \textbf{imbalanced}: only $36.9\%$ of pairs are duplicates versus $63.1\%$
non-duplicates, so accuracy alone is misleading and we report \textbf{F1} and
\textbf{log loss} as primary metrics. Second, the \emph{word-overlap ratio}
between the two questions is by far the strongest single separating signal,
which motivates the hand-crafted features described below. Questions are short,
averaging about $11$ words and mostly under $20$. The three missing question
fields are filled with the empty string.

\paragraph{Shared split for reproducibility.}
The corpus is divided into \textbf{323{,}432} training, \textbf{40{,}429}
validation, and \textbf{40{,}429} test examples (an $80/10/10$ stratified split,
\texttt{random\_state}\,$=42$). We freeze the split \emph{indices} and share them
across the team via version control, so every model family (traditional ML,
Siamese LSTM, Bi-/Cross-Encoder) is trained and evaluated on exactly the same
partitions. This makes the cross-model comparison in Section~3 a
fair, like-for-like benchmark. To prevent leakage, the TF-IDF vocabulary and all
token dictionaries are \textbf{fit on the training split only}.

\paragraph{Model-specific cleaning.}
Rather than one canonical preprocessing, we load the raw text once and branch
into three pipelines whose cleaning intensity is matched to the consumer model
(Table~\ref{tab:pipelines}). Traditional ML and the LSTM receive aggressively
normalized text (lowercasing, contraction expansion, punctuation removal);
numbers are kept for ML (e.g.\ ``iphone 14'') but dropped for the LSTM, whose
embedding vocabulary contains few numerals. BERT instead receives \emph{minimal}
cleaning (HTML and whitespace only), because its WordPiece tokenizer handles
casing and punctuation internally.

\begin{table}[h]
    \centering
    \caption{Three preprocessing pipelines branched from a single load.}
    \label{tab:pipelines}
    \begin{tabular}{lll}
        \toprule
        \textbf{Pipeline} & \textbf{Target model} & \textbf{Cleaning} \\
        \midrule
        \texttt{for\_ml}   & Traditional ML  & lowercase + expand contractions + strip punct., keep numbers \\
        \texttt{for\_lstm} & Siamese LSTM    & same as above, but remove numbers \\
        \texttt{for\_bert} & BERT Cross-Enc. & minimal (HTML \& whitespace only) \\
        \bottomrule
    \end{tabular}
\end{table}

\paragraph{Features for traditional ML.}
The \texttt{for\_ml} pipeline produces two complementary representations. (i) A
set of \textbf{10 hand-crafted features} capturing length and lexical-overlap
structure: per-question character length and word count, absolute length
difference and mean length, common-word and total-word counts, the
\textbf{word-share} ratio $|A\cap B|/(|A|+|B|)$, and the cosine similarity
between the two TF-IDF vectors. (ii) A \textbf{TF-IDF} representation
(unigrams, English stopwords removed, \texttt{min\_df}\,$=2$,
\texttt{sublinear\_tf}), where each question is vectorized separately and the two
vectors are horizontally stacked so a linear model can attend to each question
independently, yielding a sparse $323{,}432 \times 93{,}012$ training matrix.

% ----------------------------------------------------------------------------
\subsection{Traditional ML Baselines}
% ----------------------------------------------------------------------------

The traditional baselines establish how far purely lexical signals carry on this
task and serve as the reference point for the neural models. We compare
\textbf{eight} classifiers across the two feature views above: Logistic
Regression (on hand features and on TF-IDF), a calibrated Linear SVM, Multinomial
Naive Bayes, Decision Tree, Random Forest, Gradient Boosting, and XGBoost.

\paragraph{Feature engineering: \texttt{tfidf\_word\_share}.}
For the tree ensembles we add an IDF-weighted version of word overlap inspired by
the \emph{anokas} Kaggle baseline. Whereas plain word-share counts every shared
word equally, this feature weights each shared word by its inverse document
frequency,
\[
  \texttt{tfidf\_word\_share}
  = \frac{\sum_{w \in A\cap B} \mathrm{idf}(w)}
         {\sum_{w \in A\cup B} \mathrm{idf}(w)},
\]
so that overlap on \emph{rare, content-bearing} words counts for more than
overlap on common words. It is computed directly from the stored TF-IDF matrices
(no re-tokenization), bringing the ensemble hand-feature count to $11$.

\paragraph{Tuning protocol.}
Hyperparameters are selected by \textbf{minimizing validation log loss} over a
small grid per model (e.g.\ the regularization strength $C$ for the linear models,
\texttt{alpha} for Naive Bayes, and \texttt{max\_depth}/\texttt{learning\_rate}/
\texttt{subsample} for the ensembles). We tune against the held-out validation
split directly rather than with $k$-fold cross-validation, because the
$323$K-row training set makes repeated CV refits prohibitively expensive; the
frozen split keeps this selection honest. Test metrics in Table~\ref{tab:tradml}
are reported only for the refit best configuration of each model.

\begin{table}[h]
    \centering
    \caption{Traditional ML baselines on the shared test split. Best value per
    F1 / log-loss column---our imbalance-aware selection metrics---in
    \textbf{bold}. Boosted trees on hand features win on F1 and log loss; the
    TF-IDF linear models only lead on raw accuracy.}
    \label{tab:tradml}
    \small
    \begin{tabular}{llccccc}
        \toprule
        \textbf{Model} & \textbf{Features} & \textbf{Val F1} & \textbf{Test Acc} & \textbf{Test F1} & \textbf{Log loss} & \textbf{Size (KB)} \\
        \midrule
        Logistic Regression & Hand    & 0.499  & 0.655 & 0.505 & 0.563 & 0.8 \\
        Logistic Regression & TF-IDF  & 0.630  & 0.755 & 0.628 & 0.512 & 727 \\
        Linear SVM          & TF-IDF  & 0.621  & 0.746 & 0.640 & 0.514 & 2{,}182 \\
        Naive Bayes         & TF-IDF  & 0.636  & 0.745 & 0.633 & 0.599 & 2{,}907 \\
        Decision Tree       & Hand    & 0.628  & 0.709 & 0.632 & 0.508 & 40 \\
        Random Forest       & Hand    & 0.647  & 0.727 & 0.646 & 0.487 & 414{,}826 \\
        Gradient Boosting   & Hand    & 0.644  & 0.729 & 0.645 & 0.484 & 1{,}766 \\
        \textbf{XGBoost}    & Hand+   & \textbf{0.653} & 0.733 & \textbf{0.651} & \textbf{0.480} & 4{,}188 \\
        \bottomrule
    \end{tabular}
\end{table}

\paragraph{Findings.}
Because the classes are imbalanced, we rank models by \textbf{F1 and log loss}
rather than accuracy. By that criterion the hand-feature boosted trees win:
\textbf{XGBoost is the strongest baseline overall}, with the best F1 ($0.651$)
and the lowest, best-calibrated log loss ($0.480$), because its IDF-weighted
overlap and length features recover more of the minority duplicate class. The
TF-IDF linear models post slightly higher raw accuracy (LR $0.755$) but lower F1
($\sim$0.62--0.64); accuracy mostly rewards predicting the majority
non-duplicate class, so under the imbalance-aware metric these models fall behind
the ensembles. \texttt{tfidf\_word\_share} is consistently the most important
feature in the ensembles. Notably, the best tree models reach this F1 with a
$\sim$4\,MB checkpoint and sub-millisecond inference, whereas Random Forest pays
a $400$\,MB footprint for slightly worse numbers.

\paragraph{Limitation $\rightarrow$ motivation for deep models.}
Across every classifier and feature set, F1 plateaus around $0.65$
(and accuracy in the $73$--$76\%$ range). The ceiling is structural: all of these
features ultimately measure \emph{lexical overlap}, so they cannot recognize a
duplicate that is phrased with different words. For example, ``How to learn
Python'' and ``Best way to study Python'' share almost no tokens yet express the
same intent, and every lexical model here misses it. Closing this semantic gap is
precisely what motivates the contextual models that follow---the Word2Vec\,+\,LSTM
encoder and the Transformer Bi-/Cross-Encoders---whose results in
Section~3 confirm that modeling meaning, rather than word overlap,
is what lifts F1 past the traditional ceiling.
